\documentclass[11pt,a4paper]{article}
\usepackage[margin=25mm]{geometry}
\usepackage{amsmath,amssymb,amsthm}
\newtheorem{theorem}{Theorem}
\title{The regular polytopes of the third Conway group}
\author{Human Compression project: L071 recovery}
\date{Working manuscript; separate scientific review pending}
\begin{document}\maketitle
We retain the complete proofs of enumeration by stabilizers of initial
strings and of its implementation through double cosets in centralizers.
Isomorphisms preserve the order of the distinguished generators. In
particular, a polytope and its dual are counted separately unless they
are isomorphic. This is the convention of Leemans--Mixer,
\emph{Algorithms for classifying regular polytopes with a fixed
automorphism group}, Contributions to Discrete Mathematics 7(2)
(2012), 105--118, Section~7.

\begin{theorem}
The third Conway group is the automorphism group of exactly
$22\,908$ regular polytopes up to isomorphism. Their rank distribution is
\[
\begin{array}{c|rrr}
\text{rank}&3&4&5\\\hline
\text{number}&21\,118&1\,746&44.
\end{array}
\]
Every other rank has count zero. The complete distribution by ordered
Schl\"afli type agrees with every entry of the tables in Sections~7.1--7.3
of the source, including every printed zero and every unlisted type.
The exact comparison data comprise $225$ printed rank-three cells,
$565$ rank-four cells and $34$ rank-five cells.
\end{theorem}
\begin{proof}
Use the degree-$276$ permutation model of $G=\mathrm{Co}_3$ in
GAP's primitive-group library. Its identification with $\mathrm{Co}_3$
is a retained library premise. The execution records its generators
literally and checks
\[
 |G|=495\,766\,656\,000,\qquad Z(G)=1,
 \qquad |\operatorname{Aut}(G)|=|G|.
\]
Consequently conjugation by $G$ realizes every automorphism of $G$;
the compatible action required by the retained enumeration theorem
is the natural action of $G$ itself. The completed enumeration also
checks the exact order of the conjugation image against the order
of the computed full automorphism group.

For an accepted prefix $t=(t_1,\ldots,t_m)$, its automorphism
stabilizer is $S=C_G(t_1,\ldots,t_m)$ and the group containing
the eligible next involutions is
$Q=C_G(t_1,\ldots,t_{m-1})$. The retained double-coset lemma
selects exactly one $S$-orbit from every involution conjugacy class
in $Q$. The recursive end-intersection criterion then accepts
exactly the extensions with the intersection property. Its premise
is the already accepted prefix; it is not applied to arbitrary
tuples without that premise.

The completion test uses the normal closure in $Q$ of all the
eligible class representatives. This is precisely the subgroup
generated by every eligible involution, rather than just the selected
representatives. Every later mark lies in that subgroup, so the
test discards only prefixes that cannot generate $G$. The complete
retained induction proves that the generating leaves represent
each automorphism orbit once.

Every new mark lies outside the subgroup generated by its predecessors.
The subgroup order therefore increases by a factor of at least two.
Since
\[
                2^{38}\le |G|<2^{39},
\]
the search through rank $38$ covers every possible marking. The
completed execution records all leaf tuples and all zero rows
through this bound. It uses neither the source counts nor a
source-imposed rank-five cutoff to direct or stop the search.
Its numbers of generating leaves are the three displayed counts.
There are no leaves at any other rank, including ranks one and two.

A separate execution checks all $22\,908$ recorded tuples. For
each tuple it verifies the orders of the marks, their membership
in the literal group, generation, every string commutation and
the full intersection property for every pair of subsets of the
marks. For a tuple $(t_1,\ldots,t_r)$ it computes the ordered type
\[
              (|t_1t_2|,\ldots,|t_{r-1}t_r|).
\]
Grouping those computed types gives the complete distribution.
All $824$ printed source cells agree, and every computed nonzero
type is listed by the source. Thus the comparison also checks the
implicit zero assertion for every unlisted type. The source
transcription and its physical pages have been separately inspected.
The verification of recorded objects supplements the completeness
proof; it does not replace that proof or its completed execution.
Finally the retained correspondence between regular polytopes and
string C-groups identifies these ordered automorphism orbits with
the asserted isomorphism classes of polytopes.
\end{proof}

\paragraph{Attribution and finite premises.}
This recovers the source classification; it is not a new classification
claim. Stabilizer enumeration is the method described by Leemans--Mixer,
there attributed to Hartley--Hulpke. Double cosets, normal closure and
the end-intersection criterion are explicitly retained in the complete
common proofs. Exact GAP group operations, the complete conjugacy and
automorphism algorithms, and the named permutation model remain
software and library premises. The full literal outputs are retained.
There is no claim of Lean verification.

The complete L071 recovery also requires its other finite classifications,
compatible-action criterion and explicitly stated qualifications. This
component alone adds no whole-original or compression-factor credit.
\end{document}
